\documentclass[10pt]{article}
\usepackage[margin=0.72in]{geometry}
\usepackage[T1]{fontenc}
\usepackage{lmodern}
\usepackage{microtype}
\usepackage{amsmath,amssymb}
\usepackage{booktabs}
\usepackage{xcolor}
\usepackage{hyperref}
\hypersetup{
  colorlinks=true,
  urlcolor=blue!55!black,
  linkcolor=blue!55!black,
  pdftitle={Issue 128 D5-Integrated Machine Certificate},
  pdfauthor={Junkai Wang (WangTheoPhys)}
}
\setlength{\parindent}{0pt}
\setlength{\parskip}{0.34em}
\newcommand{\code}[1]{\texttt{#1}}

\title{\textbf{Issue 128: D5-Integrated Machine Certificate}\\
\large A certified $4.1357418449\times$ resource reduction}
\author{Junkai Wang (WangTheoPhys)}
\date{31 July 2026}

\begin{document}
\maketitle

\begin{abstract}
For the unchanged periodic $12\times12$ spin-$\tfrac12$ isotropic Heisenberg
benchmark at $T=1$ and global operator-norm tolerance $10^{-6}$, a new
post-freeze certificate accepts 95 fourth-order Suzuki steps and rejects 94.
The published rigorous control uses 393 steps.  Under the same compilation
model, merged group exponentials decrease from 11,791 to 2,851, giving the
exact ratio $11791/2851=4.135741844966678\ldots$.  The improvement integrates
an exact, SHA-256-bound degree-five Pauli sidecar into the finite-step error
ledger.  Fast verification, deep coefficient regeneration, adjacent-step
minimality, and corruption-rejection tests all pass.  No fivefold result is
claimed.
\end{abstract}

\section{Pinned comparison and certified resources}

The Hamiltonian, $L=12$, $N=144$, periodic boundary conditions, time, error
metric, four-matching decomposition, 31-stage five-copy fourth-order Suzuki
formula, and compiled-cost rules are identical on both sides.  Adjacent stage
merging gives $G(r)=30r+1$ group exponentials.  Each matching contains 72
bonds and the common upper-bound model charges three CNOTs per bond
propagator.

\begin{table}[h]
\centering
\begin{tabular}{@{}lrrr@{}}
\toprule
Resource & Published control & D5-integrated & Reduction factor \\
\midrule
Trotter steps & 393 & 95 & $393/95$ \\
Merged group exponentials & 11,791 & 2,851 & $11791/2851$ \\
Bond propagators & 848,952 & 205,272 & $11791/2851$ \\
CNOT upper bound & 2,546,856 & 615,816 & $11791/2851$ \\
\bottomrule
\end{tabular}
\caption{Integer resource accounting under one unchanged cost model.}
\end{table}

Thus the primary resource claim is
\[
 \frac{G(393)}{G(95)}=\frac{11791}{2851}
 =4.13574184496667835847\ldots .
\]
Relative to the control this removes 643,680 bond propagators and an upper
bound of 1,931,040 CNOTs.

\section{Proof refinement}

The frozen schema-v3 result used a verified anticommuting D4 partition but a
generic D5 majorant.  The new certificate keeps the D4 proof unchanged and
replaces only that D5 majorant with an exact grouped sidecar.  It contains
605,832 canonical Pauli coefficients partitioned into 123,106 same-support,
pairwise-anticommuting groups of size at most ten.  Its exact site-norm upper
bound is
\[
 \frac{44948270001027856175670154896253}
 {4000000000000000000000000000000}.
\]

For each group $K_g=\sum_j c_jP_j$, pairwise anticommutation permits the
outward rational bound
\[
 \lVert K_g\rVert\le
 \sqrt{\sum_j \operatorname{upper}(|c_j|)^2}.
\]
The grouping search is untrusted.  The verifier independently checks payload
canonicality and SHA-256, exact coverage and absence of duplicates, support
equality, every symplectic anticommutation relation, every outward square-root
bound, and all metadata copied into the main certificate.

\section{Finite-step ledger and adjacent boundary}

\begin{table}[h]
\centering
\begin{tabular}{@{}lr@{}}
\toprule
Right-generator contribution at $r=95$ & Global upper bound \\
\midrule
D4 & $5.7218800760132647\times10^{-7}$ \\
D5, exact grouped & $3.4853484615103630\times10^{-8}$ \\
D6 & $1.7893140877453180\times10^{-7}$ \\
D7 & $2.1921361303553973\times10^{-8}$ \\
D8 and above, geometric tail & $1.8255487798872918\times10^{-7}$ \\
\midrule
Total & $\mathbf{9.9044914028324506\times10^{-7}}$ \\
\bottomrule
\end{tabular}
\caption{Outward-readable decimals projected from exact rational values.}
\end{table}

All acceptance decisions use exact rational arithmetic, not the displayed
decimals.  At the adjacent step $r=94$, the same function evaluates to
$1.0468061165603706\times10^{-6}$, above the tolerance.  Hence 95 is the
smallest integer accepted by this certificate.  Fast verification now
regenerates each displayed contribution and the adjacent-step value.  In
particular, setting D6, D7, or the tail to zero and synchronously recomputing
the submitted total is rejected.

\section{Verification and reproducibility}

Fast verification is suitable for a live demonstration:
\begin{verbatim}
PYTHONPATH=src python3 scripts/verify.py \
  certificates/issue128-d5-integrated-certificate.json
\end{verbatim}
It returns \code{valid=true}, \code{candidate\_steps=95},
\code{candidate\_group\_exponentials=2851},
\code{d5\_term\_count=605832}, and the exact ratio
\code{11791/2851}.  It also returns
\code{finite\_step\_bound\_recomputed=true}.  Deep verification is invoked
with \code{--deep}.  It scans all 31 published-theorem centers, selects center
20, regenerates all D4 and D5 coefficients from the Suzuki formula, and
requires exact interval equality with the two sidecars.  The current-source
deep run took 408.88 seconds, used at most 1,538,768,896 bytes of resident
memory, and returned \code{deep\_proof\_regenerated=true} and
\code{baseline\_centers\_scanned=31}.

A second verifier uses only the Python standard library:
\begin{verbatim}
python3 scripts/reference_verify.py \
  certificates/issue128-d5-integrated-certificate.json
\end{verbatim}
It independently checks sidecar coverage, symplectic anticommutation,
square-root domination, and equality of three D4/D5 metadata values: term
count, group count, and maximum group size.  It also checks every finite-step
contribution, the Suzuki tail, the 95/94 comparison, and resource arithmetic
in 3.63 seconds.  It does not regenerate the coefficient maps; that separate
obligation belongs to deep mode.  The normal suite reports 146 passed tests
and 12 explicitly deselected slow tests; ten focused certificate mutations
and eight independent-checker tests pass.

On a nondegenerate open $2\times3$ diagnostic with seven unique bonds, dense
evaluation at 95 steps gives operator-norm error
$1.8313133889571672\times10^{-10}$, below the outward scaled periodic
comparison $4.129166666666667\times10^{-8}$.  This checks normalization and
stage ordering; it is not a proof that the periodic certificate applies to
open boundaries.

The original ten-file delivery manifest remains frozen at 97 steps and
$11791/2911$.  It has not been silently rewritten.  The new artifact
\code{issue128-d5-integrated-certificate.json} is separately reproducible from
the old certificate and the existing D5 sidecar using
\code{scripts/build\_d5\_integrated\_certificate.py}.

\section{Fivefold boundary}

Fivefold resource arithmetic requires $r\le78$.  However, the current D4
contribution alone is approximately $1.259084\times10^{-6}$ at $r=78$, already
larger than the entire tolerance.  Nonnegative D5--tail contributions can only
increase the ledger.  Exact-D8 HPC work therefore remains diagnostic: it may
guide a future proof, but it cannot establish fivefold without a new D4
tightening.  This report makes no fivefold claim.

\section{Claim boundary and release gates}

This certificate is scoped to the stated Hamiltonian, normalization, periodic
boundary, formula, tolerance, and cost model.  It does not prove that 94
physical steps fail, that the grouping heuristic is optimal, or that the
method is state-of-the-art over unrestricted quantum simulation algorithms.

A release is accepted only when:
\begin{itemize}
  \item the 95-step primary and standard-library reference verifiers pass;
  \item deep mode reports 31 scanned centers and coefficient regeneration;
  \item all ordinary and adversarial mutation tests pass;
  \item the D5-integrated SHA-256 manifest is clean; and
  \item this PDF is rebuilt from the committed \LaTeX{} source and visually
        inspected page by page.
\end{itemize}

\end{document}
